%% LyX 2.5.1 created this file.  For more info, see https://www.lyx.org/.
%% Do not edit unless you really know what you are doing.
\documentclass[brazil]{article}
\usepackage[T1]{fontenc}
\usepackage[utf8]{luainputenc}
\usepackage{geometry}
\geometry{verbose,tmargin=3cm,bmargin=3cm,lmargin=3cm,rmargin=2.5cm}
\usepackage{babel}
\begin{document}
\title{Trabalhos relacionados}

\maketitle
Conforme mencionado anteriormente, a ideia de trabalhar com aprendizado de máquina em computadores quânticos já é bem conhecida dentro da academia, sendo assim é de se esperar a existência de outros trabalhos na área, nos quais este próprio trabalho também se baseou.

O primeiro artigo que se destaca neste aspecto é um artigo de 2019 de Tacchino e colaboradores. Neste artigo, foi proposto e implementado um modelo híbrido de treinamento para um perceptron binário, além disso também são propostos 2 algoritmos (bloco de mudança de sinal de HSGS) para o cálculo do produto interno no computador quântico e é sugerida a possibilidade da construção de um modelo contínuo utilizando as fases dos números complexos envolvidos \cite{Italiano}. Desta forma, este artigo serviu como base para a implementação do modelo híbrido de treinamento para o perceptron quântico, assim como inspiração para a elaboração do modelo contínuo.

Outro artigo mais recente, datado de 2019 é de Havlíček e colaboradores \cite{havlivcek2019supervised}. Diferente do artigo anterior, que é baseado na proposta e implementação de rede neurais, este propõe e implementa um modelo baseado em método kernel, também baseado em um esquema de treinamento híbrido. Assim, esta pesquisa serviu como inspiração para a escolha do terceiro modelo proposto. Porém, o modelo original apresentado pelo artigo se aproxima do estado da arte e requer elevado grau de domínio da área, o que exigiria um tempo de estudo além do disponível para a elaboração do presente trabalho. Por esta razão, optou-se por implementar um modelo simplificado como introdução ao método kernel em computadores quânticos.

Em 2020, Gabor e colaboradores publicaram um interessante artigo sobre a inteligência artificial quântica, no qual é discutido o estado da arte e os principais desafios atuais da área \cite{gabor2020holy}. O autores classificam como o ``Santo Graal da IA Quântica'' o problema do \textit{feedback}, isto é, substituir os \textit{loops} que usualmente ocorrem durante o treinamento de diferentes modelos de inteligência artificial por um algoritmo inteiramente quântico. Porém, ao reconhecer a necessidade de processar uma quantidade relativamente grande de dados de treinamento, pondera-se sobre a possibilidade de que isto pode impedir que o \textit{loop} de treinamento seja inteiramente removido. Então é proposto que a chave para um maior desenvolvimento da IAQ (Inteligência Artificial Quântica) está na combinação e hibridização de vários algoritmos e técnicas, isto é, utilizar abordagens híbridas que permitam a construção de etapas de pré/pós processamento clássico.

Em outro artigo também de 2020, Lloyd e colaboradores propõem um modelo híbrido utilizando o método kernel através de uma estratégia que chamam de métrica quântica de aprendizagem \cite{lloyd2020quantum}. Esta estratégia poderia ser implementada em plataformas de computação quântica de médio prazo, apesar de incluir a representação de dados clássicos como estados quânticos no espaço de Hilbert via uma função denominada mapa de recursos quânticos (\textit{quantum feature map}), processo conhecido como incorporação quântica (\textit{quantum embedding}). Este é um processo que em geral não pode ser reproduzido em computadores clássicos, portanto deixa em aberto a possibilidade de uma vantagem quântica. É interessante lembrar outro artigo mais antigo no qual Lloyd fez parte, datado de 2014 \cite{Rebentrost_2014}. Foi um estudo liderado por Rebentrost que foi apoiado dentre outros pela DARPA (Agência de Projetos de Pesquisa Avançada de Defesa dos EUA) e pelo Laboratório de Inteligência Artificial Google-NASA. Neste artigo é proposto um modelo também baseado no método kernel, considerado um algoritmo quântico para trabalhar com grandes quantidades de dados (\textit{big data}), e tem como característica de destaque maior segurança na privacidade dos dados envolvidos no processo quando comparado com alternativas clássicas, além de um ganho de velocidade no processamento.

Winci e colaboradores publicaram um artigo também em 2020 no qual destacam que o desenvolvimento de modelos híbridos quânticos-clássicos são o estado da arte da área. Neste trabalho foi proposto também um modelo híbrido de redes neurais no qual sugerem que pode ser o caminho em direção de alcançar uma vantagem quântica, isto é, através do seu modelo o computador quântico seria capaz de resolver problemas no qual a computação clássica se mostrou lenta e ineficiente \cite{Winci_2020}. Do mesmo ano e também discutindo um possível caminho para a vantagem quântica, temos um artigo do Coyle colaboradores \cite{vantagem}.

Porém, antes de discutir qual seria este caminho, este trabalho nos apresenta a definição do que seria a chamada supremacia quântica: a capacidade de dispositivos quânticos fornecerem soluções eficientes para problemas que não podem ser resolvidos em um tempo polinomial por meios clássicos. De acordo com Coyle, a demonstração desta capacidade seria uma demonstração da chamada supremacia quântica. Neste artigo, além de propor um modelo híbrido de aprendizado de máquina, os autores também propõem a existência de uma chamada supremacia quântica de aprendizagem, isto é, uma vantagem de dispositivos e modelos quânticos sobre suas alternativas clássicas dentro do contexto de aprendizagem de máquina. Esta é uma proposta baseada na possibilidade de que existam algumas classes de distribuições que não podem ser eficientemente aprendidas por computadores clássicos, mas que podem ser aprendidas em dispositivos quânticos.

Por fim, Zhang e Ni, também em 2020, publicaram um artigo no qual discutem os avanços recentes no aprendizado de máquina quântico, nos dando uma visão geral do estado da arte. Sendo um artigo de revisão, ele possui uma estrutura bastante interessante, inclusive para um primeiro contato com a área, pois nos seus dois primeiros tópicos apresenta alguns conceitos de computação quântica para aqueles que estão familiarizados apenas com a mecânica quântica, além de discutir alguns algoritmos já conhecidos da computação quântica, como o Algoritmo de Grover e o HHL \cite{ultimo}. Na sequência, discute a aplicação de alguns algoritmos para o aprendizado de máquina quântico. É interessante aqui destacar o Support Vector Machine (SVM), que é um método kernel e aparece repetidamente em diversos artigos e as redes neurais quânticas, que foram os dois modelos no qual este TCC se baseou. Além disso, é discutido sobre o algoritmo de clusterização conhecido como K-means quântico e alguns algoritmos quânticos de redução de dimensionalidade (QPCA e QLDA).

Concluindo, os autores nos lembram que apesar de muitos pesquisadores terem demonstrado um ganho no tempo de processamento nas versões quânticas dos algoritmos de máquina de aprendizado em comparação com suas contra-partidas clássicas, ainda há alguns problemas na área. Como por exemplo a entrada e saída de dados.

\bibliographystyle{plain}
\bibliography{referencias}

\end{document}
